\verseline{}{与此同时，埃涅阿斯已率领船队循着}
\verseline{}{既定航程，在北风中劈开幽黑的波浪，}
\verseline{}{不断回望城墙；那里，不幸的厄利莎的}
\verseline{}{烈焰已经照亮全城。究竟是什么点起如此大火，}
\verseline{5}{起因仍然隐秘；可是伟大爱情受污后的酷烈痛苦}
\verseline{}{以及狂怒中的女子可能做出的事，他们深知}
\verseline{}{这些念头化作凶兆，牵引特乌克罗斯人的心。\footnote{城墙的火光来自狄多的葬火，远去的特洛伊人尚不知起因。他们想到深情受辱后的痛苦，以及狂怒女子可能做出的事，因而感到不祥。}}
\verseline{}{船只驶入深海，再没有一片}
\verseline{}{陆地迎面出现，四周唯有海水与天空；}
\verseline{10}{一团深蓝的雨云停在他们头顶，}
\verseline{}{带来黑夜和风暴，波浪在幽暗中耸起寒栗。}
\verseline{}{舵手帕利努鲁斯本人也从高高的船尾喊道：}
\verseline{}{“唉！为什么这样浓重的云团包围了天空？}
\verseline{}{父神涅普顿啊，你又准备做什么？”他说完，}
\verseline{15}{便命人收拢帆具，伏身在坚实的船桨上，}
\verseline{}{又把鼓起的帆斜转迎风，开口说道：}
\verseline{}{“宽宏的埃涅阿斯啊，即使朱庇特亲自担保，}
\verseline{}{这样的天色也不能让我指望抵达意大利。}
\verseline{}{风向已经改变，从幽黑的西方横扫而来，}
\verseline{20}{呼啸着一齐涌起，空气凝聚成云。}
\verseline{}{我们既无力逆风，也再不能沿着原定的}
\verseline{}{航向前进。既然时势胜过了我们，就顺从它，}
\verseline{}{它召唤去哪里，我们便转向哪里。我想，}
\verseline{}{兄弟厄律克斯那可靠的海岸和西坎尼亚港口并不遥远}
\verseline{25}{只要我还准确记得从前观测过的星辰。”}
\verseline{}{虔敬的埃涅阿斯答道：“我早已看出风在这样要求，}
\verseline{}{也看见你徒然同它抗争。}
\verseline{}{转帆改道吧。还有哪片土地能让我更加亲近}
\verseline{}{还有哪里更值得我停下疲惫的船只}
\verseline{30}{能胜过守护达耳达诺斯人阿刻斯忒斯}
\verseline{}{又把父亲安喀塞斯遗骨抱在怀中的土地？”}
\verseline{}{话音一落，他们便驶向港口；顺遂的西风}
\verseline{}{吹满船帆，舰队飞快地越过波涛，}
\verseline{}{终于欣然转向那片熟悉的沙岸。}
\verseline{35}{远处，阿刻斯忒斯从高高的山巅惊见}
\verseline{}{盟友的船只归来，便赶来迎接；}
\verseline{}{他携着森然标枪，披利比亚母熊皮，粗犷威武；}
\verseline{}{一位特洛伊母亲从克里尼索斯河神受孕}
\verseline{}{将他生下。他不忘古老的父祖}
\verseline{40}{欣然祝贺流亡者归来，用乡野的丰饶}
\verseline{}{接待他们，以亲厚的馈赠抚慰疲惫的友人。\footnote{阿刻斯忒斯是西西里英雄，由特洛伊女子与克里尼索斯河神所生。他接待埃涅阿斯一行，后来收容留居者，治理新城。}}
\verseline{}{次日，明亮的白昼刚从东方驱散群星，}
\verseline{}{埃涅阿斯便从整片海岸召来同伴集会，}
\verseline{}{站在坟丘的高处说道：}
\verseline{45}{“伟大的达耳达诺斯子孙，源自天神崇高血脉的族人，}
\verseline{}{十二个月已经走完，周年又至：}
\verseline{}{正是在一年前，我们把遗骸与神圣父亲的骨殖}
\verseline{}{安置入土，奉献了哀悼的祭坛。}
\verseline{}{如今，若我没有记错，那一天已经来临；我将永远}
\verseline{50}{以苦痛、也永远以敬意纪念它——众神如此安排。}
\verseline{}{即使我流亡在盖图利亚的叙耳提斯沙洲，}
\verseline{}{或被困在阿尔戈斯海上、迈锡尼城中，}
\verseline{}{仍会践行周年的誓愿，依礼举行庄严的队列，}
\verseline{}{把各自相宜的礼物堆上祭坛。}
\verseline{55}{如今，我们竟不经谋划便来到父亲的骨灰与遗骨前；}
\verseline{}{在我看来，众神的心意与神力必在其中：}
\verseline{}{我们被送来，进入这友好的港口。}
\verseline{}{来吧，让我们一同欣然举行这场祭礼；}
\verseline{}{祈求顺风，也愿父亲允许我在建城之后，}
\verseline{60}{年年把这些圣礼带到为他奉献的神庙。}
\verseline{}{特洛伊所生的阿刻斯忒斯按船只数目，}
\verseline{}{每船赠你们两头牛；请把祖传的家神，}
\verseline{}{也把主人阿刻斯忒斯所敬奉的神祇请来赴宴。}
\verseline{}{此外，倘若第九个黎明为凡人带来温煦的}
\verseline{65}{白昼，以光芒重新揭开大地，}
\verseline{}{我将先为特乌克罗斯人设下快船的竞逐；}
\verseline{}{无论谁善于疾足奔跑，谁勇力过人，}
\verseline{}{能以标枪和轻箭胜过众人，}
\verseline{}{或敢凭生皮拳具交手，}
\verseline{70}{都请前来，等待胜者应得的奖赏。\footnote{埃涅阿斯纪念安喀塞斯死去一周年，并希望建城后每年在献给父亲的神庙举行祭礼；此次祭赛定于第九个黎明。塞尔维乌斯（5.64）以罗马的第九日葬仪解释这一安排。葬礼竞技承接《伊利亚特》第二十三卷阿喀琉斯为帕特罗克洛斯举行竞技的传统。}}
\verseline{}{人人噤声敬神，用枝叶环绕额头。”}
\verseline{}{他说完，以母亲所爱的桃金娘遮住鬓边；}
\verseline{}{赫吕穆斯这样做，年事已高的阿刻斯忒斯这样做，}
\verseline{}{少年阿斯卡尼乌斯也这样做，其余青年随之效仿。}
\verseline{75}{埃涅阿斯离开集会，率领成千上万的人}
\verseline{}{走向坟丘，庞大的人群簇拥在他周围。}
\verseline{}{他在那里依礼奠下两杯纯酒，献给巴克科斯，}
\verseline{}{又向地面浇下两杯鲜奶、两杯祭牲的鲜血，}
\verseline{}{撒上紫色花朵，开口说道：}
\verseline{80}{“致敬，神圣的父亲；再次致敬啊，从特洛伊救出}
\verseline{}{却终归徒然的骨灰，以及父亲的灵魂和亡影。}
\verseline{}{命运不许我同你一起寻找意大利的疆界、命定的田野，}
\verseline{}{以及奥索尼亚的台伯河——无论它身在何处。”}
\verseline{}{话音刚落，一条巨蛇便从坟墓深处滑出，}
\verseline{85}{以七道巨环、七重盘曲的身躯，}
\verseline{}{温和地环抱坟丘，从祭坛间游过；}
\verseline{}{深蓝的斑纹布满它的背，金色的点点光辉}
\verseline{}{燃亮鳞片，仿佛云中的彩虹}
\verseline{}{迎着太阳投下千种变幻的颜色。}
\verseline{90}{埃涅阿斯看得惊住了。巨蛇拖着长长的身躯，}
\verseline{}{终于游到祭酒盘与光洁的杯盏之间，}
\verseline{}{尝过祭食，又无害地退回}
\verseline{}{坟丘深处，离开它享用过的祭坛。}
\verseline{}{埃涅阿斯因而更加郑重地重行父亲的祭礼；}
\verseline{95}{他不知道该把它视为此地之神，还是父亲的}
\verseline{}{侍从，便依照旧俗宰杀两只祭羊，}
\verseline{}{两头猪和同样数目的黑背小公牛，}
\verseline{}{从祭酒盘倾下葡萄酒，呼唤}
\verseline{}{伟大的安喀塞斯之魂，以及从阿刻戎获释而来的亡灵。\footnote{蛇与坟墓、亡灵和地方守护神常有联系。埃涅阿斯不确定它是此地的守护神，还是父亲的侍从；incertus 明写他的迟疑。}}
\verseline{100}{同伴们也欣然按照各自所有}
\verseline{}{带来礼物，堆满祭坛，宰杀小公牛；}
\verseline{}{有些人依次架起青铜大锅，有些人散坐草地，}
\verseline{}{把炭火放在烤叉下，炙烤祭牲的内脏。}
\verseline{}{期待已久的日子到了；法厄同的马已经用光明}
\verseline{105}{载来晴朗的第九个黎明。}
\verseline{}{传闻与著名的阿刻斯忒斯之名召来了邻近居民；}
\verseline{}{他们欢欣聚集，挤满海岸，}
\verseline{}{有的来观看埃涅阿斯的族人，有的准备亲自竞逐。}
\verseline{}{礼物先被陈列在众人眼前，放在场地}
\verseline{110}{中央：神圣的鼎、青翠的花冠、}
\verseline{}{胜者所得的棕榈枝、兵器、浸染紫红的}
\verseline{}{衣袍，以及数塔兰同的金银；}
\verseline{}{号角从中央土丘吹响，宣告祭赛开始。\footnote{这里的 ludi、certamina 是安喀塞斯周年祭中的祭赛。竞赛开始前，三足鼎、花冠、棕榈枝、兵器、紫染衣物以及金银已经陈列在场中；各项结束后，埃涅阿斯再按赛果颁发相应奖品。}}
\verseline{}{四艘从整支舰队选出的船只势均力敌，}
\verseline{115}{以沉重的船桨开始了第一场竞逐。}
\verseline{}{门涅斯透斯率骁勇的桨手驾驶迅捷的“普里斯提斯号”——}
\verseline{}{后来的意大利人门涅斯透斯，梅弥乌斯氏族之名由他而来；}
\verseline{}{古阿斯驾驶庞然如城的“喀迈拉号”，}
\verseline{}{达耳达诺斯青年分作三层推动巨舰，}
\verseline{120}{三列船桨同时起落；}
\verseline{}{塞耳革斯图斯乘着巨大的“半人马号”——塞耳吉乌斯氏族由他得名；}
\verseline{}{克洛安图斯则驾驶深蓝的“斯库拉号”，}
\verseline{}{罗马的克卢恩提乌斯氏族啊，你们正是他的后裔。\footnote{门涅斯透斯、塞耳革斯图斯和克洛安图斯分别被追溯为罗马梅弥乌斯、塞耳吉乌斯与克卢恩提乌斯氏族的祖先。诗歌越过故事时代，把奥古斯都时代的罗马家族接入特洛伊祭赛，也由此为这些氏族赋予神话祖源。}}
\verseline{}{远处海上有一块礁石，正对着浪花飞溅的}
\verseline{125}{海岸；每当冬季西北风遮没星辰，}
\verseline{}{它便沉入涨起的海浪，任波涛捶打。}
\verseline{}{风平浪静时，它默默露出水面，像一片平坦场地}
\verseline{}{立在不动的波上，是海鸟最爱的向阳栖处。}
\verseline{}{父亲埃涅阿斯在那里竖起一根取自枝叶繁茂的圣栎、}
\verseline{130}{仍带绿叶的标柱，给水手作记号，使他们知道}
\verseline{}{从哪里返航，在哪里转过漫长的航程。}
\verseline{}{他们随后抽签选定位置；船长们站在船尾，}
\verseline{}{黄金与紫衣装点周身，从远处灿然生辉；}
\verseline{}{其余青年以白杨叶覆住额头，}
\verseline{135}{赤裸的肩膀浸过油，闪着光泽。}
\verseline{}{他们坐上桨座，手臂绷紧在桨柄上；}
\verseline{}{凝神等候信号；怦然的恐惧攫住}
\verseline{}{跃动的心，求取荣誉的渴望高高扬起。}
\verseline{}{清亮的号角一响，所有船只立刻}
\verseline{140}{从起点跃出；水手的呼声撞上天空，}
\verseline{}{臂膀用力回收，翻开的海水涌起白沫。}
\verseline{}{它们同时犁开水道；整个海面}
\verseline{}{被船桨和三齿船首撕裂，豁然张开。}
\verseline{}{双马战车从起栅放出、}
\verseline{145}{奔腾着扑向原野，也不曾这般迅猛；}
\verseline{}{车手纵放套马，也不曾这样抖动}
\verseline{}{飘荡的缰绳，俯身悬在挥落的鞭子上。}
\verseline{}{人群的拍掌、喧声与偏爱者的热望}
\verseline{}{顿时使整片树林一齐轰鸣；闭合的海岸}
\verseline{150}{卷回呼喊，受声浪撞击的山丘处处应和。}
\verseline{}{古阿斯在众船前冲出，在人群的}
\verseline{}{喧腾中滑上最先的波头；克洛安图斯跟在后面，}
\verseline{}{桨手更强，却受沉重的松木船身拖慢。}
\verseline{}{在他们之后，“普里斯提斯号”和“半人马号”相隔无几，}
\verseline{155}{都想夺取前方的位置；}
\verseline{}{时而“普里斯提斯号”领先，时而巨大的“半人马号”}
\verseline{}{越过对手，时而两船并行，船首齐平，}
\verseline{}{长长的龙骨一同犁过咸水。}
\verseline{}{船只已经接近礁石，即将抵达标柱；}
\verseline{160}{领先的古阿斯正在海心占尽优势，}
\verseline{}{忽然高声斥责舵手墨诺提斯：}
\verseline{}{“为什么往我右边偏得那么远？向这里修正航向；}
\verseline{}{贴住岸边，让左桨只擦过礁石；}
\verseline{}{让别人走外海。”他说完，墨诺提斯却害怕}
\verseline{165}{看不见的暗礁，把船头扭向海上的波浪。}
\verseline{}{“你要偏到哪里去？”古阿斯再次高喊，}
\verseline{}{“驶向礁石，墨诺提斯！”他正大声召回舵手，}
\verseline{}{回头却看见克洛安图斯已经逼到船尾，占住更近的内道。}
\verseline{}{克洛安图斯从古阿斯的船和轰响的礁石之间}
\verseline{170}{擦过左侧水道，由内抢进，忽然超过}
\verseline{}{领先者，绕过标柱，进入安全的水面。}
\verseline{}{巨大的痛楚顿时在青年骨髓里燃烧，}
\verseline{}{泪水也没有放过他的脸颊；他忘了自己的体面，}
\verseline{}{也忘了同伴的安危，把迟钝的墨诺提斯}
\verseline{175}{从高高的船尾一把掀进海里。}
\verseline{}{他亲自接过舵柄，自己做舵手和船长，}
\verseline{}{催促船员，把舵转向海岸。}
\verseline{}{年老的墨诺提斯一身湿衣，身子沉重，}
\verseline{}{终于勉强从深水中重新浮起，}
\verseline{180}{爬上礁石顶端，坐在干燥的岩面上。}
\verseline{}{特乌克罗斯人看见他坠水，又奋力游泳，不禁发笑；}
\verseline{}{见他从胸中吐出咸涩的海水，更笑个不停。}
\verseline{}{此时，落在末尾的两人心中燃起快意的希望——}
\verseline{}{塞耳革斯图斯与门涅斯透斯都想超过受阻的古阿斯。}
\verseline{185}{塞耳革斯图斯抢在前面，靠近礁石；}
\verseline{}{然而领先的还不到整个船身，}
\verseline{}{只有一部分在前，竞逐的“普里斯提斯号”以船首紧逼它的后部。}
\verseline{}{门涅斯透斯从船身中央穿过同伴之间，}
\verseline{}{开口激励他们：“现在，现在挺身压住船桨，}
\verseline{190}{赫克托耳的战友们，我在特洛伊命运的最后关头}
\verseline{}{选为同伴的人们；现在使出你们}
\verseline{}{曾在盖图利亚叙耳提斯、伊奥尼亚海，}
\verseline{}{以及马雷阿紧追不放的波浪中用过的力量与勇气。}
\verseline{}{我门涅斯透斯已不求第一，也不再非胜不可；}
\verseline{195}{虽然——唉！但涅普顿啊，你把胜利赐给谁，就让谁胜出吧。}
\verseline{}{只是莫让我们最后归来；同胞们，胜过这件耻事，}
\verseline{}{不许如此不堪的结果发生！”他们竭尽全力，}
\verseline{}{俯身压桨；青铜包覆的船身因重击而震颤，}
\verseline{}{水面从船底向后退去；急促的喘息摇撼四肢}
\verseline{200}{和干渴的嘴唇，汗水从全身流淌如溪。}
\verseline{}{恰是一场意外，为他们带来期盼的荣誉。}
\verseline{}{塞耳革斯图斯心神狂热，把船首挤向}
\verseline{}{礁石内侧，闯入狭窄而不利的水道，}
\verseline{}{不幸搁在突出的岩石上。}
\verseline{205}{礁石震动；撑在尖锐石脊上的船桨}
\verseline{}{喀然折断，撞坏的船首悬挂在那里。}
\verseline{}{水手们一齐起身，高声叫喊，被迫停下；}
\verseline{}{他们取出包铁的撑杆和尖头长篙，}
\verseline{}{又从水中收拢折断的船桨。}
\verseline{210}{门涅斯透斯却因这番转机而欢喜，锐气倍增；}
\verseline{}{船桨迅疾而有节律地起落，他又呼来风力，}
\verseline{}{冲下海面，在开阔水域疾驰。}
\verseline{}{仿佛一只栖居岩穴的鸽子忽然受惊，}
\verseline{}{那多孔岩石里藏着它的家与可爱的巢；}
\verseline{215}{它飞向田野，惊惶的翅膀在洞穴}
\verseline{}{拍出巨响，旋即滑入安静的天空，}
\verseline{}{轻掠澄澈的道路，迅疾的双翼不再振动：}
\verseline{}{门涅斯透斯正是这样，“普里斯提斯号”正是这样切过}
\verseline{}{最后一段水面，奔行的冲力载着它飞翔。}
\verseline{220}{它先把塞耳革斯图斯抛在高高的礁石上挣扎；}
\verseline{}{他陷在浅水里，徒然呼求}
\verseline{}{援助，正学习怎样用断桨前进。}
\verseline{}{门涅斯透斯接着追上古阿斯和庞大的}
\verseline{}{“喀迈拉号”；它失去了舵手，只得退让。}
\verseline{225}{终点就在眼前，如今只剩克洛安图斯领先；}
\verseline{}{他奋尽全力追向对方，步步逼迫。}
\verseline{}{呼声骤然加倍；所有人都以各自的热望}
\verseline{}{催促追赶者，喧响震得天空回荡。}
\verseline{}{领先的人若守不住本有的荣誉和已经赢得的光荣，}
\verseline{230}{便深以为耻，宁愿拿生命交换声名；}
\verseline{}{追赶者却从顺遂中获得力量：他们能做到，因为自信能够做到。}
\verseline{}{两只船本可能船首齐平地取得奖品，}
\verseline{}{若不是克洛安图斯向海面伸出双手，}
\verseline{}{倾吐祷词，呼唤众神接受誓愿：}
\verseline{235}{“统治大海、让我驶过你们水域的神祇啊，}
\verseline{}{我将欣然在这片海岸的祭坛前，为你们}
\verseline{}{安置一头雪白的公牛，履行我的誓约；我要把内脏}
\verseline{}{献入咸涩的波浪，倾下清澈的葡萄酒。”}
\verseline{}{他说完，深层波涛之下所有涅瑞伊得斯、}
\verseline{240}{福耳库斯之众和少女帕诺佩娅都听见了；}
\verseline{}{港口之神波耳图努斯更亲自伸出巨手，}
\verseline{}{推动前行的船；它比南风和飞箭更快，}
\verseline{}{直奔陆地，藏进深深的港湾。\footnote{克洛安图斯仍居领先，门涅斯透斯正在逼近；他向海神许献白牛、内脏和葡萄酒。众海神听见祈祷，波耳图努斯亲手推船，助他先抵港口。波耳图努斯是罗马的港口神。}}
\verseline{}{安喀塞斯之子随即依照惯例召来众人，}
\verseline{245}{由传令官高声宣布克洛安图斯获胜，}
\verseline{}{又以青翠月桂环绕他的额头；}
\verseline{}{他准许各船选取三头小公牛作为礼物，}
\verseline{}{还赠予葡萄酒和一大塔兰同的白银。}
\verseline{}{他又为每位船长另添特别的荣誉：}
\verseline{250}{给胜者一件绣金斗篷，四周有一道又一道}
\verseline{}{墨利波亚紫纹，曲折成双重回环；}
\verseline{}{其上织着一位王族少年，在枝叶繁茂的伊达山}
\verseline{}{投掷标枪，以奔跑追逐迅捷的雄鹿，}
\verseline{}{英姿勃发，仿佛仍在喘息；疾飞的朱庇特神鹰}
\verseline{255}{以弯曲的利爪将他从伊达山攫向高空；}
\verseline{}{年迈的护卫徒然伸掌指向星辰，}
\verseline{}{猎犬的狂吠在空中激荡。\footnote{斗篷织着特洛伊王族少年伽倪墨得斯被劫的故事：他在伊达山狩猎，朱庇特的神鹰将他攫往天上，成为众神的斟酒者；织物上，老护卫徒然伸手，猎犬朝天空狂吠。}}
\verseline{}{给凭本领取得第二名的船长，}
\verseline{}{他赠送一件以光滑环扣缀成、三重金丝交织的}
\verseline{260}{铠甲；那是他从前在高耸的伊利乌姆城下、湍急的西摩伊斯河边}
\verseline{}{得胜后亲自从得摩勒俄斯身上剥下来的，}
\verseline{}{如今归这位男子所有，在战阵中既是荣耀也是护卫。}
\verseline{}{两个仆人斐革乌斯和萨伽里斯肩膀一同用力，}
\verseline{}{才勉强抬动这件层叠的重甲；得摩勒俄斯当年}
\verseline{265}{却曾身披它奔跑，驱赶四散的特洛伊人。}
\verseline{}{第三份礼物是两个青铜大釜，}
\verseline{}{还有一对精制银碗，表面浮雕凸起。}
\verseline{}{众人已经领过礼物，为所得的财富自豪，}
\verseline{}{以紫红束带环绕额头，正要离去；}
\verseline{270}{忽然，塞耳革斯图斯费尽手段，才挣脱严酷的礁石，}
\verseline{}{失去了船桨，一侧桨列也已残缺，}
\verseline{}{驾着受人嘲笑、没有荣誉的船驶来。}
\verseline{}{它仿佛一条在高起的道路上被困住的蛇，}
\verseline{}{青铜车轮斜斜碾过它，或有行人猛然抡起}
\verseline{275}{石块重击，把它残破而半死地遗在路上；}
\verseline{}{它想逃走，长长的身体却徒然扭曲，}
\verseline{}{前半身依旧凶猛，双眼燃烧，昂起}
\verseline{}{嘶鸣的颈项；受伤而瘸折的后半身却把它拖住，}
\verseline{}{它挣扎着盘成圈结，又把身体反折回自身：}
\verseline{280}{那艘迟缓的船就靠这样残缺的桨列移动；}
\verseline{}{但它终于升起船帆，满帆驶入港口。}
\verseline{}{埃涅阿斯因塞耳革斯图斯保住船只、带回同伴而欣喜，}
\verseline{}{便把许诺的礼物给了他。}
\verseline{}{他得到一名女奴，熟谙弥涅耳瓦所司的织作，}
\verseline{285}{名叫福洛厄，来自克里特，怀里还有一对吃奶的孪生儿子。\footnote{福洛厄是克里特女奴，与尚在哺乳的孪生儿子一同作为赠礼，体现史诗中奴隶及其子女可作为财产的处境。弥涅耳瓦主管织作，诗以女神之名称她所熟悉的技艺。}}
\verseline{}{这场竞逐结束后，虔敬的埃涅阿斯走向}
\verseline{}{一片草地；四周弯曲的山丘}
\verseline{}{和树林环绕着它，谷地中央形成天然剧场，}
\verseline{}{设有环形场地。埃涅阿斯率领成千上万的人}
\verseline{290}{来到坐席中央，在搭起的高台上落座。}
\verseline{}{谁愿以飞快的脚步较量，}
\verseline{}{他便以奖品激励其心，陈列赏赐。}
\verseline{}{特乌克罗斯人和西坎尼亚人从四面会聚；}
\verseline{}{首先是尼苏斯和欧律阿卢斯：}
\verseline{295}{欧律阿卢斯以出众的容貌和蓬勃青春引人注目，}
\verseline{}{尼苏斯则因对少年的忠贞深情而引人注目；随后是}
\verseline{}{出自普里阿摩斯尊贵王族的狄俄瑞斯；}
\verseline{}{萨利乌斯和帕特戎一同跟来，前者是阿卡纳尼亚人，}
\verseline{}{后者出自阿卡狄亚忒革亚的血统；}
\verseline{300}{还有两位特里纳克里亚青年赫吕穆斯和帕诺佩斯，}
\verseline{}{惯于山林，是年迈阿刻斯忒斯的随从；}
\verseline{}{此外还有许多人，声名幽微，没有留下他们的名字。\footnote{尼苏斯对欧律阿卢斯的 amor pius 是忠贞的深情。赛跑中，尼苏斯以阻挡萨利乌斯帮助欧律阿卢斯取胜；第九卷两人还会共同夜袭，并一同死去。}}
\verseline{}{埃涅阿斯站在他们中间，开口说道：}
\verseline{}{“请用心听取这些话，让心情欢畅：}
\verseline{305}{你们之中没有一人会不受赠礼便离开。}
\verseline{}{我将给每人两支克诺索斯标枪，}
\verseline{}{铁尖磨得明亮，再给一把嵌银的雕花双刃斧；}
\verseline{}{这是人人共有的荣誉。最先的三人另有奖赏，}
\verseline{}{都将以金黄的橄榄枝环绕额头。}
\verseline{310}{第一位胜者可得一匹饰有华美马具的骏马；}
\verseline{}{第二位可得一只亚马逊箭袋，装满色雷斯}
\verseline{}{箭矢，宽阔的金色背带环绕在外，}
\verseline{}{以一枚圆润宝石制成的搭扣系紧；}
\verseline{}{第三位便带着这顶阿尔戈斯头盔离去。”}
\verseline{315}{话音一落，他们各就各位；信号蓦然响起，}
\verseline{}{他们立刻飞奔过赛道，离开起点，}
\verseline{}{如骤雨般倾泻，一同奔向终点。}
\verseline{}{尼苏斯最先奔出，远远跃在所有人之前，}
\verseline{}{比风、比闪电的羽翼还快；}
\verseline{320}{萨利乌斯跟在后面，虽说最近，}
\verseline{}{两人之间却仍相隔甚远；再落下一段，}
\verseline{}{第三位是欧律阿卢斯；}
\verseline{}{赫吕穆斯紧追着他；在赫吕穆斯身后，}
\verseline{}{看哪，狄俄瑞斯已经飞到脚跟贴着脚跟，}
\verseline{325}{俯近他的肩头；如果再多一段路程，}
\verseline{}{他或许已抢到前面，至少会使胜负难分。}
\verseline{}{他们几乎已经跑完赛程，疲惫地}
\verseline{}{接近终点；不幸的尼苏斯忽在滑腻的血上失足——}
\verseline{}{那里先前恰好宰杀过小公牛，鲜血流在地面，}
\verseline{330}{浸湿了青翠的草叶。}
\verseline{}{青年已经为胜利欢欣，脚踏在湿滑的地面}
\verseline{}{却没能稳住踉跄的步子，向前扑倒}
\verseline{}{摔进污秽的粪泥与神圣的祭血。}
\verseline{}{即使如此，他仍没有忘记欧律阿卢斯，没有忘记自己的爱；}
\verseline{335}{他从滑地上起身，故意挡在萨利乌斯面前；}
\verseline{}{萨利乌斯翻滚着倒在厚密的沙中。}
\verseline{}{欧律阿卢斯冲了出去，借朋友的帮助夺得}
\verseline{}{第一，在赞许的拍掌和欢声中飞过终处。}
\verseline{}{随后是赫吕穆斯；第三枚棕榈枝如今归狄俄瑞斯。}
\verseline{340}{萨利乌斯高声抗议，叫喊响彻巨大剧场的全部坐席}
\verseline{}{也传到前排长者面前；}
\verseline{}{他要求归还被诡计夺走的荣誉。}
\verseline{}{但众人的偏爱护住欧律阿卢斯，他惹人怜爱的泪水也护住他；}
\verseline{}{美德从俊美身体显现，总显得更加可爱。}
\verseline{345}{狄俄瑞斯也高声出言相助；}
\verseline{}{他已升到得奖的位置，若把第一的荣誉}
\verseline{}{归还萨利乌斯，自己便徒然来到最后一份奖赏前。}
\verseline{}{父亲埃涅阿斯于是说道：“你们各自的礼物都确定不移；}
\verseline{}{这个少年的棕榈枝也无人会从原位挪动；\footnote{Conte 2019 校勘记认可 C. W. Nauck 将 pueri 与 palmam 相连的属格解释：所指为欧律阿卢斯，“这个少年的棕榈枝”保持原来的获奖次序。}}
\verseline{350}{但请准许我怜惜无辜朋友所遭的不幸。”}
\verseline{}{他说完，把一张巨大的盖图利亚狮皮}
\verseline{}{赠给萨利乌斯；浓密鬃毛压在上面，爪尖闪着金色。}
\verseline{}{尼苏斯这时说道：“如果落败者得到的礼物如此丰厚，}
\verseline{}{你又怜悯跌倒的人，那么什么礼物}
\verseline{355}{才配得上尼苏斯？我本来凭本领应获第一的花冠，}
\verseline{}{若非敌意的运气像掀翻萨利乌斯那样掀翻了我？”}
\verseline{}{他边说边展示沾满湿泥的面庞，}
\verseline{}{以及污损得可笑的肢体。慈厚的父亲对他笑了，}
\verseline{}{命人取来一面狄狄马翁精工制作的盾牌，}
\verseline{360}{达那俄斯人曾从涅普顿神庙的圣柱上把它拆下。}
\verseline{}{他以这份卓越的赠礼赏赐了出众的青年。}
\verseline{}{赛跑结束，礼物分发完毕，他又说道：}
\verseline{}{“现在，谁若胸中有勇力与坚定的心，}
\verseline{}{就请上前，缠好双掌，把手臂举起。”}
\verseline{365}{他说完，为拳斗设下两份奖赏：}
\verseline{}{胜者可得一头以黄金和圣带装饰的公牛，}
\verseline{}{败者也有一柄剑和一顶华美的头盔作为慰藉。}
\verseline{}{没有迟疑；达瑞斯立刻凭借强大的力量}
\verseline{}{昂首站出，在众人巨大的低鸣中挺身而立。}
\verseline{370}{从前只有他惯于同帕里斯较量；}
\verseline{}{也正是他，在伟大的赫克托耳安息的坟墓旁，}
\verseline{}{将身躯巨大的胜者布忒斯——这人}
\verseline{}{自称来自柏布律喀亚的阿米库斯一族——}
\verseline{}{击倒，使他垂死躺在黄沙上。}
\verseline{375}{达瑞斯就以这样的气势首先高昂着头，走向搏斗，}
\verseline{}{显出宽阔的双肩；他伸展手臂，}
\verseline{}{左右挥动，一拳拳击打空气。}
\verseline{}{人们为他寻找对手；偌大人群中却没有一人}
\verseline{}{敢于上前，敢把皮革拳具缠上双手迎战。}
\verseline{380}{他于是兴冲冲地走来，以为所有人都已让出胜利，}
\verseline{}{站在埃涅阿斯脚前；不再耽搁，}
\verseline{}{左手握住公牛的角，这样说道：}
\verseline{}{“女神之子，既然没有人敢把自己交给这场搏斗，}
\verseline{}{我还要站到几时？还该把我留在这里多久？}
\verseline{385}{下令让我领走礼物吧。”达耳达诺斯子孙一齐高呼，}
\verseline{}{要把许诺的奖赏交给他。}
\verseline{}{严肃的阿刻斯忒斯这时出言责备恩忒卢斯；}
\verseline{}{恩忒卢斯恰坐在他旁边一片青翠的草垫上：}
\verseline{}{“恩忒卢斯，你从前最勇于众人，难道只是徒有虚名？}
\verseline{390}{你竟能这样忍耐，任由如此丰厚的}
\verseline{}{礼物未经搏斗便被拿走？你总提起那位神圣的师父}
\verseline{}{厄律克斯，如今他在哪里？你传遍}
\verseline{}{特里纳克里亚的名声，还有挂在屋顶的战利品，又在哪里？”}
\verseline{}{恩忒卢斯答道：“我对声名和荣耀的热爱并未}
\verseline{395}{被恐惧赶走；只是衰老使冰冷的}
\verseline{}{血液迟钝，身体里衰竭的力量也已寒凉。}
\verseline{}{如果我仍有从前那种青春，}
\verseline{}{仍有这个放肆的家伙仗恃而得意的年华，}
\verseline{}{根本无须奖赏和漂亮的公牛引诱，}
\verseline{400}{我就会走上前来；我并不在意礼物。”说完，}
\verseline{}{他把一对重得惊人的皮革拳具}
\verseline{}{抛在场地中央；刚烈的厄律克斯从前惯于}
\verseline{}{带着它们交手，把坚硬的牛皮缠上双臂。}
\verseline{}{众人惊住了：那巨大的拳具由七层}
\verseline{405}{牛皮叠成，里面缝入铅块和铁块，僵硬无比。}
\verseline{}{达瑞斯本人最为骇然，远远退开，拒绝使用；}
\verseline{}{宽宏的安喀塞斯之子也把那分量和}
\verseline{}{一层层巨大的皮带翻来覆去细看。}
\verseline{}{老人于是从胸中发出这样的声音：}
\verseline{410}{“倘若有人见过赫尔库勒斯本人的拳具和兵器，}
\verseline{}{见过就在这片海岸发生的惨烈搏斗，又会怎样？}
\verseline{}{你的亲兄弟厄律克斯从前就戴着这副拳具——}
\verseline{}{你看，至今还染着鲜血与飞溅的脑浆——}
\verseline{}{他曾凭它迎战伟大的阿尔喀得斯；我也曾惯于使用，}
\verseline{415}{那时更好的血液还给我力量，嫉妒的老年尚未}
\verseline{}{把白色零星撒上两边鬓角。}
\verseline{}{可是，特洛伊人达瑞斯若拒绝用我们的拳具，}
\verseline{}{虔敬的埃涅阿斯若也同意，发起祭赛的阿刻斯忒斯若赞成，}
\verseline{}{那就让搏斗公平。我收起厄律克斯的牛皮拳具，}
\verseline{420}{别再害怕；你也脱下特洛伊人的拳具。”\footnote{Caestus 是以皮带缠在手掌与前臂上的古代拳具。厄律克斯旧拳具由七层牛皮叠成，又缝入铅铁，具有致命危险；达瑞斯拒绝使用后，埃涅阿斯为双方换上分量相等的拳具，才使搏斗继续。}}
\verseline{}{说完，他从肩头掀下双层斗篷，}
\verseline{}{裸露出庞大的肢体、巨骨与上臂，}
\verseline{}{高大的身躯站到沙地中央。}
\verseline{}{父亲安喀塞斯之子于是取来分量相等的拳具，}
\verseline{425}{把同样的装备缠上两人的手掌。}
\verseline{}{两人立刻踮起脚尖，挺身站定，}
\verseline{}{毫无惧色地把手臂举向高空。}
\verseline{}{他们高昂的头远远后仰，躲开拳锋，}
\verseline{}{手与手相交，开始寻衅进击。}
\verseline{430}{一人脚步更灵活，又有青春可恃；}
\verseline{}{另一人强在体格和重量，但迟缓的膝盖}
\verseline{}{颤抖不稳，艰难的喘息摇撼着庞大四肢。}
\verseline{}{两人徒然挥出许多拳，又一次次}
\verseline{}{重击彼此内凹的肋部，胸膛发出巨响；}
\verseline{435}{拳头频频游走在耳边和鬓侧，}
\verseline{}{颌骨在坚硬的打击下喀喀作响。}
\verseline{}{沉重的恩忒卢斯站稳不动，始终维持同一姿势，}
\verseline{}{只轻移身体，凭警觉的双眼避开拳锋。}
\verseline{}{达瑞斯却像一个以攻城器械进逼高城的人，}
\verseline{440}{又像披甲围困山间堡垒的军队，}
\verseline{}{一会儿试探这里，一会儿试探那里，以种种技巧}
\verseline{}{走遍每一条入口，又不断变换突击，始终徒劳地逼迫。}
\verseline{}{恩忒卢斯挺身亮出右拳，高高}
\verseline{}{举起；达瑞斯看见这一击从头顶落来，}
\verseline{445}{敏捷地抽身闪开。}
\verseline{}{恩忒卢斯把全部力量倾入空中，}
\verseline{}{沉重的身躯反被巨大的分量带着}
\verseline{}{重重倒地，仿佛厄律曼托斯山或巍峨的}
\verseline{}{伊达山上，一棵根部拔起的空心松树轰然倒下。}
\verseline{450}{特乌克罗斯人和特里纳克里亚青年因各自的偏爱一齐起身；}
\verseline{}{喊声直冲天空，阿刻斯忒斯第一个奔来，}
\verseline{}{怜惜地从地上扶起同龄的友人。}
\verseline{}{但这场意外没有拖慢英雄，也没有把他吓住；}
\verseline{}{他更猛烈地重返搏斗，愤怒唤醒了力量。}
\verseline{455}{羞耻与自知的勇力燃起他的气力；}
\verseline{}{他烈火般追赶达瑞斯，使他在整片场地仓皇退却，}
\verseline{}{时而右拳连击，时而换作左拳；}
\verseline{}{没有停顿，没有歇息：乌云降下的密集冰雹}
\verseline{}{怎样在屋顶喀喀作响，这位斗者也怎样}
\verseline{460}{以两手的重拳频频捶击，把达瑞斯打得旋转不定。}
\verseline{}{父亲埃涅阿斯不许怒气继续发展，}
\verseline{}{不许恩忒卢斯带着苦烈的心肠逞凶；}
\verseline{}{他终止搏斗，把疲惫的达瑞斯}
\verseline{}{救出来，以话语安慰他，这样说道：}
\verseline{465}{“不幸的人，何等失常竟攫住了你的心？}
\verseline{}{你还没有感到另一种力量，感到神意已经转向吗？}
\verseline{}{向神明退让吧。”他说完，便以声音分开搏斗。}
\verseline{}{达瑞斯忠实的同伴把他搀向船只；他拖着无力的膝盖，}
\verseline{}{头颅向两边摇晃，口中吐出浓稠的血，}
\verseline{470}{牙齿混杂在血里。}
\verseline{}{他们领命接过头盔和剑，}
\verseline{}{把棕榈枝与公牛留给恩忒卢斯。}
\verseline{}{胜者意气高昂，又以公牛为荣：}
\verseline{}{“女神之子，还有你们特乌克罗斯人，”他说，“请由此看清}
\verseline{475}{我年轻的身体里曾经有过怎样的力量，}
\verseline{}{也看清你们从怎样的死亡中召回了达瑞斯。”}
\verseline{}{他说完，便迎着公牛站定；}
\verseline{}{那公牛立在那里，充作拳斗的奖赏。他把右手向后收起，}
\verseline{}{挺身把硬皮拳具瞄准两角之间，}
\verseline{480}{猛然砸下，击碎骨骼与脑髓。}
\verseline{}{公牛倒地，失去生命，颤抖着伏在土上。}
\verseline{}{恩忒卢斯在它上方从胸中倾吐这些话：}
\verseline{}{“厄律克斯，我把这条更好的生命献给你，代替达瑞斯之死；}
\verseline{}{我已获胜，就在这里放下拳具，也放下这门技艺。”\footnote{恩忒卢斯把本可落在达瑞斯身上的致命一击转向公牛，并称它为献给厄律克斯、代替达瑞斯之死的“更好的生命”。击牛既展示他的力量，也是一场明确的替代性献祭；他随即放下拳具与技艺，宣布退出搏斗。}}
\verseline{485}{埃涅阿斯随即邀请愿以迅箭}
\verseline{}{较量的人前来，并宣布奖赏；}
\verseline{}{他以巨手从塞雷斯图斯的船上竖起一根桅杆，}
\verseline{}{把一只扑翼的鸽子系在穿过桅顶的绳上，}
\verseline{}{高高悬起，作为铁箭所向的目标。}
\verseline{490}{众人聚来；投入的签由一顶青铜}
\verseline{}{头盔承接。喝彩声中，首先抽到位置的}
\verseline{}{是许耳塔库斯之子希波科翁；}
\verseline{}{接着是刚在船赛中获誉的门涅斯透斯，}
\verseline{}{他仍以青翠的橄榄枝环绕额头。}
\verseline{495}{最著名的潘达鲁斯啊，第三位是你的兄弟}
\verseline{}{欧律提翁；从前正是你受命打破盟约，}
\verseline{}{首先把箭射进阿凯亚人中央。\footnote{潘达鲁斯是《伊利亚特》第四卷破坏希腊人与特洛伊人停战盟约的弓手；他受弥涅耳瓦诱导，向墨涅拉俄斯放箭。此处欧律提翁是他的兄弟。}}
\verseline{}{最后，阿刻斯忒斯的签留在头盔底部，}
\verseline{}{他也敢亲手尝试青年的劳作。}
\verseline{500}{射手们以强劲的力量弯下弓身，}
\verseline{}{各自从箭袋中取出箭矢。}
\verseline{}{弓弦发出颤响，第一支箭穿过天空，}
\verseline{}{许耳塔库斯之子的飞箭劈开疾驰的空气；}
\verseline{}{它直飞而去，扎进迎面的桅杆木身。}
\verseline{505}{桅杆震颤，受惊的鸟拍打翅膀，}
\verseline{}{四下都响起它振翅的巨声。}
\verseline{}{刚烈的门涅斯透斯随后拉弓站定，}
\verseline{}{瞄准高处，同时把目光和箭矢指向目标。}
\verseline{}{可惜他不能以铁锋碰到鸟身，}
\verseline{510}{却射断了绳结与亚麻的束缚；}
\verseline{}{鸽子的脚正被它们系在高高的桅杆上，}
\verseline{}{那鸟飞入南风，逃进乌云。}
\verseline{}{欧律提翁动作迅捷，早已握着}
\verseline{}{搭在拉满弓上的箭，以誓愿呼求兄长；}
\verseline{515}{他望见鸽子已在空旷的天空欢快飞翔，}
\verseline{}{便射中那在黑云下拍动双翼的鸟。}
\verseline{}{鸽子失去生命，坠落下来；它把性命留在}
\verseline{}{高处的星空，又把扎在身上的箭带回地面。}
\verseline{}{只剩阿刻斯忒斯尚未出手，胜者的棕榈枝却已失去；}
\verseline{520}{他仍把箭射向高空，}
\verseline{}{老人仍要显示自己的技艺，也让弓弦鸣响。}
\verseline{}{忽然，一个异象出现在众人眼前，}
\verseline{}{后来将成为重大征兆；其后发生的大事才显明它，}
\verseline{}{令人畏惧的先知很久以后才唱出迟来的预示。}
\verseline{525}{那支芦箭飞翔在澄澈的云间，忽然燃烧，}
\verseline{}{以火焰标出道路，渐渐消逝，}
\verseline{}{耗尽在风中；仿佛星辰常从天空挣脱，}
\verseline{}{飞掠而过，拖出长长的光发。}
\verseline{}{特里纳克里亚与特乌克罗斯人心神震动，久久呆立，}
\verseline{530}{一齐向天神祈祷；伟大的埃涅阿斯也不拒绝}
\verseline{}{这一征兆。他拥抱欣喜的阿刻斯忒斯，}
\verseline{}{把丰厚的赠礼堆给他，说道：}
\verseline{}{“收下吧，父亲；奥林匹斯的伟大君王以这样的}
\verseline{}{兆象表示，要让你得到特别的荣誉。}
\verseline{535}{你将得到年迈安喀塞斯本人留下的这件礼物：}
\verseline{}{一只压印图纹的调酒钵，从前色雷斯人}
\verseline{}{基修斯曾把它作为厚礼送给父亲安喀塞斯，}
\verseline{}{留作基修斯的纪念，也作为情谊的信物。”}
\verseline{}{他说完，以青翠月桂环绕阿刻斯忒斯的额头，}
\verseline{540}{宣布他越过众人，成为第一位胜者。\footnote{阿刻斯忒斯射向天空的箭在云中燃烧；诗说后来的大事与先知之言将显明这个征兆，未点明具体事件。埃涅阿斯据此授予他特别荣誉，宣布他为第一位胜者，并赠以安喀塞斯传下的调酒钵。}}
\verseline{}{善良的欧律提翁并不嫉妒这份置于他之上的荣誉，}
\verseline{}{尽管只有他从高空射落了鸽子。}
\verseline{}{第二个上前领取礼物的，是射断束绳的人；}
\verseline{}{最后才是以飞箭钉中桅杆的人。}
\verseline{545}{父亲埃涅阿斯尚未遣散祭赛，}
\verseline{}{便把未成年的尤卢斯的监护人兼随从}
\verseline{}{厄皮图斯之子召到身边，附耳嘱咐：}
\verseline{}{“快去，如果阿斯卡尼乌斯已把少年的}
\verseline{}{队伍带在身边，也已排好马匹的阵势，}
\verseline{550}{就叫他为祖父率骑队上场，披甲显身。”}
\verseline{}{他说完，又亲自命令散布在长长场地上的}
\verseline{}{人群全部退开，让原野空旷。}
\verseline{}{少年们缓步上场；他们骑着勒好缰辔的马，}
\verseline{}{在父母眼前一同闪耀；看着他们走过，}
\verseline{555}{所有特里纳克里亚与特洛伊青年都惊叹低鸣。}
\verseline{}{众人依照习俗，以修整的花冠压住头发；}
\verseline{}{各持两根前端装有铁尖的山茱萸木矛；}
\verseline{}{有些人在肩上背着光滑的箭袋；胸膛最高处}
\verseline{}{环着一道柔韧的绞金圈，绕过脖颈。}
\verseline{560}{骑手共分三队，三位首领各自驰行；}
\verseline{}{每位领队身后跟着两组六名少年，}
\verseline{}{分列成阵，随同身份相称的指挥者灿然行进。}
\verseline{}{欢欣的第一列少年由小普里阿摩斯率领；}
\verseline{}{他承继祖父的名字——波利特斯啊，这是你显赫的}
\verseline{565}{后裔，日后将使意大利一族壮大——一匹色雷斯骏马}
\verseline{}{驮着他，双色毛皮缀着白斑，一只前蹄带有白色印记}
\verseline{}{高昂的额头也显出一片白斑。}
\verseline{}{第二位是阿提斯，拉丁人的阿提乌斯氏族由他而来；}
\verseline{}{小阿提斯也是少年尤卢斯所爱重的少年朋友。}
\verseline{570}{最后一位是尤卢斯，容貌胜过众人；}
\verseline{}{他骑着一匹西顿骏马，那是美丽的狄多}
\verseline{}{赠给他的，留作狄多的纪念，也是爱意的信物。}
\verseline{}{其余少年所骑的特里纳克里亚马匹}
\verseline{}{来自年迈的阿刻斯忒斯。}
\verseline{575}{达耳达诺斯人以掌声迎接惴惴不安的少年，欣然凝望，}
\verseline{}{从他们的面孔认出古老先辈的模样。}
\verseline{}{少年们欣然骑马巡过全部观众席，}
\verseline{}{让所有亲人的目光看过之后，}
\verseline{}{厄皮图斯之子从远处向就绪的少年高喊信号，甩响马鞭。}
\verseline{580}{他们步调相称地向四面散开，三列队伍}
\verseline{}{分成各自的行列；一听召唤，又重新}
\verseline{}{调转道路，举起敌对的兵器。}
\verseline{}{继而沿场地展开一次次不同的冲进和返驰，}
\verseline{}{让一个个圆环交错套叠，}
\verseline{585}{披着武装唤起战斗的幻象；}
\verseline{}{一会儿转背逃走，一会儿回矛相向，}
\verseline{}{一会儿又和平地并肩驰骋。}
\verseline{}{仿佛传说中克里特高地的迷宫，}
\verseline{}{盲目的墙壁织成道路，藏有}
\verseline{590}{千重路径的两可诡局；那无从察知、无法折返的迷途}
\verseline{}{使一切循路的标志失去作用。}
\verseline{}{特乌克罗斯人的孩子也以这样的行程}
\verseline{}{交错脚步，在演练中编织逃奔与交锋；}
\verseline{}{他们仿佛游泳的海豚，穿行湿润的海域，}
\verseline{595}{劈开卡耳帕托斯海和利比亚海。\footnote{Conte 2019 正文保留第595行的半行，不收入部分传本在 secant 后添加的 luduntque per undas（“在波浪间嬉游”）。该补文见 该版校勘记。}}
\verseline{}{阿斯卡尼乌斯日后以城墙环绕阿尔巴隆伽时，}
\verseline{}{首先重现这种行进的习俗与这些竞技，}
\verseline{}{又教导古老的拉丁人依礼举行，}
\verseline{}{一如他年少时同特洛伊少年亲自演练的方式；}
\verseline{600}{阿尔巴人再传给自己的子孙；后来伟大的}
\verseline{}{罗马承接了它，并保存这项祖传的荣誉；}
\verseline{}{如今少年演练称为“特洛伊”，队伍也叫“特洛伊骑阵”。}
\verseline{}{为神圣父亲举行的竞技到此全部完成。\footnote{少年骑阵把军事训练、游戏和祖先仪式结合起来，后来被称为 lusus Troiae，即“特洛伊游戏”或“特洛伊骑阵”，由罗马贵族少年在公共典礼中演练。}}
\verseline{}{自此，命运第一次转向，背弃了信义。}
\verseline{605}{正当众人在坟丘前以种种竞技重行庄严祭礼，}
\verseline{}{萨图尔努斯之女朱诺从天上派下伊里斯，}
\verseline{}{让风吹送她前往伊利乌姆的舰队；}
\verseline{}{女神筹谋重重，对古老的怨恨仍未餍足。}
\verseline{}{伊里斯沿千色长虹疾行，}
\verseline{610}{这位无人看见的少女从迅捷的轨道飞落。}
\verseline{}{她望见庞大的人群，环视海岸，}
\verseline{}{看见港口空荡、舰队无人看守。}
\verseline{}{远处一片孤单的海滩上，特洛伊妇女避开众人，}
\verseline{}{为失去安喀塞斯而哭泣；她们一边流泪，一边都望着}
\verseline{615}{深不可测的大海。唉，她们已经如此疲惫，前方竟还有这么多海路，}
\verseline{}{还有如此辽阔的大海——所有人发出的都是同一声哀叹。\footnote{Conte 2019 校勘记倾向把第615—616行理解为妇女怨叹的间接转述；fessis 意为“对疲惫的妇女们”。}}
\verseline{}{她们祈求一座城市，已经厌倦承受海上的劳苦。}
\verseline{}{伊里斯深谙害人之术，于是投入}
\verseline{}{她们中间，放下女神的容貌与衣饰，}
\verseline{620}{化作贝罗厄——特玛鲁斯人多律克卢斯年迈的妻子，}
\verseline{}{素有门第、声名与子女；}
\verseline{}{她以这副模样走进达耳达诺斯母亲们中央：}
\verseline{}{“可怜的人们啊，”她说，“阿凯亚人的手若早在战争中}
\verseline{}{把我们拖向死亡，让我们死在祖国城墙下，该有多好！不幸的族群}
\verseline{625}{命运还把你留下，准备交给怎样的毁灭？}
\verseline{}{特洛伊陷落以后，第七个夏天已经周转到来；\footnote{“第七个夏天”是从特洛伊陷落算起的漂泊时间标记。第一卷755行也有这一说法，全诗年月的衔接有解释分歧；本卷46—50行另说明安喀塞斯死去已满一周年。}}
\verseline{}{我们走遍一切海域与陆地，掠过无数不容人的礁岩}
\verseline{}{也循着一片片陌生的星空；我们在浩瀚大海上漂流}
\verseline{}{追逐不断逃远的意大利，被波浪翻来滚去。}
\verseline{630}{这里是兄弟厄律克斯的疆界，主人阿刻斯忒斯也在这里；}
\verseline{}{谁能阻止我们筑起城墙，给同胞一座城市？}
\verseline{}{祖国啊，从敌人手中徒然救出的家神啊，}
\verseline{}{难道再也不会有一座城墙被称为特洛伊？我在任何地方}
\verseline{}{都见不到赫克托耳的河流——克珊托斯与西摩伊斯了吗？}
\verseline{635}{来吧，跟我一同烧毁这些不祥的船只。}
\verseline{}{先知卡珊德拉的形影曾在梦中}
\verseline{}{仿佛递给我燃烧的火把：‘就在这里寻找特洛伊，}
\verseline{}{这里就是你们的家。’她这样说道。如今正该行动，}
\verseline{}{这般巨大的神兆不容迟疑。看哪，那里有四座}
\verseline{640}{涅普顿的祭坛；神明亲自把火把和勇气交给我们。”}
\verseline{}{她边说边第一个猛然抓起毁灭的火，}
\verseline{}{高举右手，竭力向远处挥动，}
\verseline{}{再把它掷出；伊利乌姆妇女的心一齐震动，惊愕失神。}
\verseline{}{人群中有一位年事最老的妇女，}
\verseline{645}{名叫皮尔戈，曾是普里阿摩斯众多王子的乳母：}
\verseline{}{“母亲们，这人不是贝罗厄，不是那位罗埃特翁女子——}
\verseline{}{多律克卢斯的妻子；请看她神圣美貌的标志}
\verseline{}{看她燃烧的双眼，看她的气息、}
\verseline{}{面容，听她的声音，再看她行走的步态！}
\verseline{650}{方才我亲自离开了贝罗厄，把她独自留在那里，}
\verseline{}{正卧病难起，又深感不平：只有她无缘这场}
\verseline{}{祭礼，不能向安喀塞斯献上应得的敬礼。”\footnote{伊里斯化作多律克卢斯之妻贝罗厄，又捏造卡珊德拉托梦传火的故事。皮尔戈从眼神、声音、气息与步态识出神性，并以真正的贝罗厄卧病为证；伪装被揭穿以后，伊里斯的神显反而把妇女的犹疑推向行动。}}
\verseline{}{她说完。}
\verseline{}{起初，母亲们仍然犹疑，以阴沉的眼光}
\verseline{655}{望着船只，徘徊在对眼前土地的苦涩眷恋}
\verseline{}{与命运召唤她们前往的王国之间；}
\verseline{}{忽然，女神举起平展的双翼飞上天空，}
\verseline{}{从云层下逸去，划出一道巨大的长虹。}
\verseline{}{她们于是被神迹震得失神，又被狂乱驱使，\footnote{妇女们原已厌倦航海、渴望城市；伊里斯的神迹又煽起狂乱。朱诺的影响退去后，她们重新认出亲人，并为焚船羞愧。}}
\verseline{660}{一齐叫喊，从内室的炉火夺取火种；}
\verseline{}{有些人剥取祭坛上的枝叶、柴条与火把，}
\verseline{}{纷纷投出。武尔坎努斯放开缰绳，烈火在}
\verseline{}{桨座、船桨和绘彩的杉木船身间狂奔。}
\verseline{}{欧墨卢斯把船只燃烧的消息传到安喀塞斯的}
\verseline{665}{坟丘与剧场的层层坐席；人们自己回头，}
\verseline{}{也看见黑色灰烬在浓烟中四散飞舞。}
\verseline{}{阿斯卡尼乌斯最先行动；他方才还欣然率领}
\verseline{}{骑阵演练，此刻便策马直奔混乱的}
\verseline{}{营地，惊惶的教官根本不能把他留住。}
\verseline{670}{“这是怎样新起的狂乱？如今又要走向哪里，}
\verseline{}{唉，可怜的同胞！你们烧的并非阿尔戈斯敌军和敌营，}
\verseline{}{却是自己的希望。看哪，我是你们的}
\verseline{}{阿斯卡尼乌斯！”他把空空的头盔掷在她们脚前；}
\verseline{}{那是他在演练中唤起战争幻象时所戴。}
\verseline{675}{埃涅阿斯也立即赶到，特乌克罗斯人的队伍随之而来。}
\verseline{}{妇女们却因恐惧向不同方向奔散，}
\verseline{}{沿海岸逃走，偷偷寻找树林或任何凹陷的}
\verseline{}{岩洞；她们厌恶自己的行动，也羞于面对天光，}
\verseline{}{心神改变，重新认出亲人，朱诺已从胸中震落。}
\verseline{680}{但烈焰并未因此放下}
\verseline{}{难以驯服的力量；麻絮在湿润的木材下仍然存活，}
\verseline{}{吐出迟缓的烟，热气一点点吞噬船身，}
\verseline{}{灾祸向整具躯体深处蔓延；}
\verseline{}{众人的力量和倾注成河的海水都毫无用处。}
\verseline{685}{虔敬的埃涅阿斯于是撕下肩头的衣袍，}
\verseline{}{伸出双掌，向众神呼求援助：}
\verseline{}{“全能的朱庇特啊，如果你还没有憎恨特洛伊人到最后一人，}
\verseline{}{如果昔日的敬义仍在某处顾念}
\verseline{}{人间的劳苦，父亲啊，请立刻使舰队逃离火焰，}
\verseline{690}{把特乌克罗斯人衰微的命运从死亡中抢救出来。}
\verseline{}{否则，倘若我罪有应得，就以敌意的雷霆}
\verseline{}{把剩余的一切送给死亡，用你的右手在此将我们掩埋。”}
\verseline{}{话音未落，一场黑色风暴便倾下前所未有的}
\verseline{}{暴雨，雷声震动大地的高处与原野；}
\verseline{695}{浑浊的骤雨从整个天空奔泻，}
\verseline{}{挟着密集的南风，带水而来，黑暗至极。}
\verseline{}{船身从上方灌满雨水，半焦的}
\verseline{}{木材全部湿透；直到热气尽灭，除去}
\verseline{}{已经毁失的四艘，所有船只都从灾祸中保全。\footnote{暴雨扑灭船火，但 quattuor amissis 明说已有四艘船毁失。船只减少，瑙忒斯随后建议让部分同伴留居西西里。}}
\verseline{700}{父亲埃涅阿斯却被这惨痛的意外震撼，}
\verseline{}{胸中反复转动巨大的忧虑，一会儿向这里，}
\verseline{}{一会儿向那里：是留居西西里的田野，}
\verseline{}{忘掉命运，还是前去取得意大利海岸？}
\verseline{}{年迈的瑙忒斯曾独受特里同尼亚的帕拉斯}
\verseline{705}{教导，女神以诸多技艺使他卓然出众；}
\verseline{}{她还向自己的先知作出神示：天神的盛怒预示什么，}
\verseline{}{或命运的秩序提出怎样的要求；}
\verseline{}{此刻，他便以这些话安慰埃涅阿斯，开口道：\footnote{这里的帕拉斯是弥涅耳瓦的称号，与第八卷的青年帕拉斯同名。Conte 2019 校勘记释706行 haec 为帕拉斯，女神向她的先知瑙忒斯传示天神盛怒与命运秩序所要求的事。709—710行转为瑙忒斯的劝告：随命运的牵引，以忍耐战胜遭际。}}
\verseline{}{“女神之子，无论命运把我们向前牵引还是向后拉回，我们都追随它；}
\verseline{710}{无论将要发生什么，一切遭际都须凭忍耐战胜。}
\verseline{}{你有神圣血统的达耳达诺斯人阿刻斯忒斯；}
\verseline{}{请接纳他加入谋议，他也愿与你相合。}
\verseline{}{把那些因船只毁失而留下的人交给他，还有那些}
\verseline{}{已经厌倦这项伟大事业、厌倦你的命运的人；}
\verseline{715}{挑出年迈的老人、疲于海洋的母亲，}
\verseline{}{以及随你而来的一切体弱者和畏惧危险的人，}
\verseline{}{允许这些疲惫者在这片土地拥有城墙；}
\verseline{}{若你准许，他们会把城市称为阿刻斯塔。”}
\verseline{}{年长友人的话燃起他的心；}
\verseline{720}{他的心神愈发被种种忧虑拉扯。}
\verseline{}{漆黑的夜乘双马车升上天空，占据天穹；}
\verseline{}{这时，父亲安喀塞斯的面容忽然仿佛从天而降，}
\verseline{}{向他倾吐出这些话语：}
\verseline{}{“儿子啊，我活着时，你曾比生命更令我珍爱；}
\verseline{725}{儿子啊，你饱经伊利乌姆命运的折磨，}
\verseline{}{我奉朱庇特之命来到这里；是他从舰队赶走火焰，}
\verseline{}{也终于从高天俯下怜悯。}
\verseline{}{听从年迈的瑙忒斯刚刚给你的绝佳建议；}
\verseline{}{把挑选出的青年、最勇敢的心灵}
\verseline{730}{带往意大利。拉提乌姆有一个强悍、习俗粗犷的族群}
\verseline{}{等你以战争征服。但在那以前，}
\verseline{}{先走近狄斯的地下宫殿，穿过阿韦尔努斯深处，}
\verseline{}{来寻找同我相会，儿子。亵渎神明的}
\verseline{}{塔耳塔洛斯与悲哀的亡影都未拘留我；我居住在虔敬者}
\verseline{735}{欢欣的聚会之所，居住在厄吕西乌姆。纯洁的西比尔}
\verseline{}{会以大量黑色祭牲的鲜血引领你到这里；}
\verseline{}{那时，你将认识自己的全部后裔，也会知道哪些城墙将赐给你。\footnote{安喀塞斯预告第六卷的冥界之行：埃涅阿斯将由库迈西比尔引导，经过阿韦尔努斯，以黑色祭牲献祭后进入地下，与身在厄吕西乌姆的父亲相会，并从他那里得知罗马后裔与未来城邦。}}
\verseline{}{现在告别吧；潮湿的黑夜正转过中途，}
\verseline{}{残酷的东方已让喘息马群的气息吹到我身上。”}
\verseline{740}{他说完，像轻烟一样逃入稀薄的空气。}
\verseline{}{埃涅阿斯喊道：“你要奔向哪里，为什么匆匆离去？}
\verseline{}{你在躲避谁，又是谁阻止你投入我的怀抱？”}
\verseline{}{他说着，便拨开灰烬，唤醒沉睡的火种，}
\verseline{}{又以虔敬的祭粉与盛满香料的香匣，}
\verseline{745}{伏身礼敬帕加马的家神和白发维斯塔的内殿。\footnote{安喀塞斯的幻影如烟消散，埃涅阿斯追问是谁阻止父亲进入自己的怀抱。这与第二卷克瑞乌萨和第六卷冥府重逢同属无法触及亡者的场景，但此处没有“三次拥抱”的完整重复。埃涅阿斯随即唤醒祭火，礼敬家神与维斯塔；与父亲亡魂的相会结束后，他以祭仪收束这一夜。}}
\verseline{}{他立即召来同伴，首先召来阿刻斯忒斯，}
\verseline{}{把朱庇特的命令、亲爱父亲的教诲，}
\verseline{}{以及心中已经确定的意向一一告诉他们。}
\verseline{}{谋议没有延迟，阿刻斯忒斯也不拒绝命令。}
\verseline{750}{他们把母亲们列入新城，把愿意留下的民众}
\verseline{}{安置在这里——这些心灵不再渴求宏伟的声名。}
\verseline{}{继续航行的人重新装设桨座，更换船上}
\verseline{}{被火焰啃噬的木材，配好船桨与索具；}
\verseline{}{人数虽少，勇力却仍为战争鲜活。}
\verseline{755}{与此同时，埃涅阿斯以犁划出城市的边界，}
\verseline{}{抽签分配房屋；他命令这城称作伊利乌姆，这片地方称作特洛伊。}
\verseline{}{特洛伊人阿刻斯忒斯为王国欢欣，}
\verseline{}{宣布听讼，召来元老，颁布法律。}
\verseline{}{随后，在厄律克斯高近星辰的峰顶，}
\verseline{760}{为伊达利乌姆的维纳斯建立神殿；安喀塞斯的坟墓}
\verseline{}{也得到一位祭司和一片广阔的圣林。\footnote{新城阿刻斯塔通常与西西里的塞革斯塔相认；埃涅阿斯以犁划界、分配住区，阿刻斯忒斯宣布听讼、召元老并颁布法律。厄律克斯山顶的维纳斯神殿属于当地维纳斯崇拜；“伊达利乌姆的”是维纳斯得自塞浦路斯伊达利乌姆圣地的称号。}}
\verseline{}{全族已经宴饮九日，在祭坛前}
\verseline{}{尽过敬礼；温和的风铺平海面，}
\verseline{}{南风再度频频吹来，召唤他们驶向深海。}
\verseline{765}{沿着弯曲的海岸，巨大的哭声陡然升起；}
\verseline{}{人们彼此拥抱，日夜拖延，不忍分离。}
\verseline{}{就连那些母亲，就连那些从前觉得}
\verseline{}{大海面目严酷、连名字都无法忍受的人，}
\verseline{}{如今也想同行，愿承受逃亡的全部劳苦。}
\verseline{770}{善良的埃涅阿斯以亲厚的言语安慰他们，}
\verseline{}{流着泪把他们托付给同族的阿刻斯忒斯。\footnote{Transcribunt 带有把成员编入新城名册的意思；留居者得到城市、房屋、法律和祭仪中心。分别时，连先前畏惧海洋的留居者也重新想随船离去。}}
\verseline{}{他又命人向厄律克斯宰献三头小公牛，向风暴诸神献一只羔羊，}
\verseline{}{然后依次解开缆绳。}
\verseline{}{他自己以修整的橄榄枝叶环绕额头，}
\verseline{775}{站在远离岸边的船首，手持祭酒盘，把内脏}
\verseline{}{抛进咸涩的波浪，倾下清澈的葡萄酒。}
\verseline{}{同伴竞相击水，以船桨扫过海面；}
\verseline{}{风从船尾升起，一路护送他们前行。\footnote{Conte 2019 将第778行排在第777行之前；该版校勘记据航行动作的先后解释为舰队先以桨驶离港口，继而受风推动。}}
\verseline{}{忧虑折磨着维纳斯；此刻她向}
\verseline{780}{涅普顿开口，从胸中倾吐怨诉：}
\verseline{}{“朱诺沉重的愤怒和永不餍足的心}
\verseline{}{迫使我，涅普顿啊，屈身百般恳求；}
\verseline{}{漫长时日不能软化她，任何虔敬也不能，}
\verseline{}{朱庇特的命令与命运都不能击溃她，使她安静。}
\verseline{785}{她以不可言说的仇恨，从佛里吉亚族群中央}
\verseline{}{吞噬那座城市还不满足；把特洛伊残众拖过}
\verseline{}{一切惩罚还不满足：她仍追逐那座被杀死的城市的}
\verseline{}{灰烬与骨殖。如此狂怒的缘由，她自己该最清楚。}
\verseline{}{你本人便是我的见证：不久前，她在利比亚水域}
\verseline{790}{突然掀起何等风暴，把整个海洋同天空}
\verseline{}{搅在一起，徒然借助埃俄卢斯的狂飙，}
\verseline{}{竟敢在你的王国这样行事。}
\verseline{}{看哪，如今她又以罪恶驱使特洛伊的母亲们，}
\verseline{}{可耻地烧毁船只，又迫使众人在舰队受损后}
\verseline{795}{把同伴留在陌生的土地。}
\verseline{}{我恳求你，准许余下的人把船帆托付给你的波浪，}
\verseline{}{平安抵达劳伦图姆的台伯河；}
\verseline{}{倘若我的祈求合乎神律，倘若命运三女神把那座城赐给他们。”}
\verseline{}{海洋深处的主宰、萨图尔努斯之子于是回答：}
\verseline{800}{“库忒瑞亚啊，你完全有神圣的权利信任我的王国；}
\verseline{}{你的出身也正源自我的海域。我同样配得这份信任：}
\verseline{}{我屡次压下天空与海洋的狂怒和巨大暴烈。}
\verseline{}{在陆地上，我对你的埃涅阿斯也同样关切——}
\verseline{}{克珊托斯与西摩伊斯可以作证。阿喀琉斯}
\verseline{805}{追赶失魂的特洛伊队列，把他们撞向城墙，}
\verseline{}{将成千上万人交给死亡；河流被尸体塞满，}
\verseline{}{发出呻吟，克珊托斯找不到道路，不能把自己}
\verseline{}{卷向海洋。那时，埃涅阿斯正同勇猛的珀琉斯之子}
\verseline{}{交锋，神明与力量都不相称；是我以中空的云雾把他}
\verseline{810}{救走，尽管我一心想从根基掀翻}
\verseline{}{那座背弃誓言、由我亲手建造的特洛伊城墙。\footnote{涅普顿所述之事见《伊利亚特》第二十卷：阿喀琉斯即将杀死埃涅阿斯时，海神以云雾救走后者，因为命运要保存达耳达诺斯的后裔。涅普顿曾因拉俄墨冬背誓而敌视特洛伊，却仍维护埃涅阿斯的既定命数。}}
\verseline{}{如今我的心意依旧不变；驱散恐惧吧。}
\verseline{}{他会如你所愿，平安抵达阿韦尔努斯的港口。}
\verseline{}{只会少一个人：他失落在漩涡之中，你将寻找他；}
\verseline{815}{一人的生命将为众人献出。”\footnote{Unum pro multis dabitur caput 的字面结构是“一颗头颅将为众人交出”。Caput 可转指一条生命，pro multis 即“一人为众人”；这句话预告帕利努鲁斯将成为那被交出的“一人”。}}
\verseline{}{父神以这些话抚平女神的胸怀，使她欢悦，}
\verseline{}{随即以黄金套住马匹，把起沫的}
\verseline{}{嚼子放入野兽口中，任手里全部缰绳奔放。}
\verseline{}{他驾着蔚蓝战车，轻捷飞过海水表面；}
\verseline{820}{波浪沉下，涨起的海面在轰鸣的车轴下}
\verseline{}{铺展平直，云团从广阔的天空逃散。}
\verseline{}{随后显出形貌各异的随从：庞大的海兽，}
\verseline{}{年迈的格劳库斯之众，伊诺之子帕莱蒙，}
\verseline{}{迅捷的特里同，以及福耳库斯的全部军旅；}
\verseline{825}{忒提斯、墨利忒和少女帕诺佩娅占据左侧，}
\verseline{}{还有尼萨厄、斯皮奥、塔利亚与库莫多刻。\footnote{涅普顿的随从包括格劳库斯之众、伊诺之子帕莱蒙、鱼尾海神特里同、福耳库斯的军旅，以及忒提斯、帕诺佩娅等海中女神。Tritones 在原文为复数，指一群特里同式海神。}}
\verseline{}{温柔的喜悦继而渗入父亲埃涅阿斯}
\verseline{}{仍悬在忧虑中的心；他命人赶快竖起全部桅杆}
\verseline{}{把船帆伸展在横桁上。}
\verseline{830}{众人同心调动帆脚索，先向左，}
\verseline{}{又向右放出帆腹；他们一齐旋转}
\verseline{}{高高的桁角，来回转向；各自的顺风载着舰队。}
\verseline{}{帕利努鲁斯走在最前，率领密集的}
\verseline{}{船列；其余船只奉命依照他的航线。}
\verseline{835}{潮湿的黑夜几乎已经抵达天空的中点；}
\verseline{}{水手们横卧在船桨下的坚硬座板上，}
\verseline{}{让四肢在平静的安眠中松弛；}
\verseline{}{这时，睡神轻盈地从以太星辰间滑落，}
\verseline{}{拨开黑暗的空气，驱散阴影，}
\verseline{840}{向你而来，帕利努鲁斯，为无辜的你带来凶险的梦；}
\verseline{}{神明坐在高高的船尾，}
\verseline{}{化作福耳巴斯的模样，从口中倾出这些话：}
\verseline{}{“伊阿索斯之子帕利努鲁斯啊，海水自己正载着舰队；}
\verseline{}{风息均匀；这一刻已经赐给休憩。}
\verseline{845}{放下头颅吧，从劳作中偷回疲惫的眼睛；}
\verseline{}{我会亲自代你片刻，承担你的职分。”}
\verseline{}{帕利努鲁斯勉强抬起双眼，对他答道：}
\verseline{}{“你叫我认不出平静海水的面孔和沉静的波浪吗？}
\verseline{}{叫我信赖这头怪物吗？我怎么能把埃涅阿斯}
\verseline{850}{托付给骗人的风？何况我已经如此多次}
\verseline{}{受晴空的假象欺骗？”}
\verseline{}{他说着，紧紧攀住舵柄，}
\verseline{}{始终不肯松手，双眼依然仰望星辰。}
\verseline{}{看哪，神明举起一根沾有勒忒水露、}
\verseline{855}{又以斯提克斯之力浸透睡意的树枝，在他两边}
\verseline{}{鬓角上摇动；他虽仍抗拒，浮游的双眼却已松开。}
\verseline{}{突如其来的睡眠刚使他的四肢开始松弛，}
\verseline{}{睡神便俯压下来，把船尾的一部分扯裂，}
\verseline{}{连同舵柄，将他头朝下抛入清澈的}
\verseline{860}{波浪；他一次又一次徒然呼唤同伴。}
\verseline{}{神明自己却振翼飞上稀薄的天空。\footnote{这里的 Somnus 是人格化的睡神。他先化作福耳巴斯劝帕利努鲁斯交出职责，遭拒后才使用沾有勒忒水露、含斯提克斯催眠之力的树枝，强迫他的眼睛闭合；帕利努鲁斯连同舵柄与一部分船尾被抛入海中。}}
\verseline{}{舰队仍然沿安全的航路驶过海面，}
\verseline{}{凭父神涅普顿的许诺，无所畏惧地前行。}
\verseline{}{他们已经被送到塞壬的礁石近旁；}
\verseline{865}{那里从前险恶，被无数人的白骨覆满，}
\verseline{}{此刻岩石在不断的海涛中从远方发出粗哑轰鸣。}
\verseline{}{父亲埃涅阿斯忽然察觉舵手失落，船只正在漂荡，}
\verseline{}{便亲自在夜间波浪上掌舵；}
\verseline{}{他一再呻吟，友人的遭际震撼了他的心：}
\verseline{870}{“啊，你太过信赖宁静的天空与海面，}
\verseline{}{帕利努鲁斯，你将赤身躺在陌生的沙岸！”\footnote{帕利努鲁斯明确拒绝相信平静海面，始终抓住舵柄仰望星辰；睡神的行动只向叙述者和读者显现，埃涅阿斯却在卷末误以为他“太过信赖”宁静的天空与海洋。第六卷将进一步说明，他坠海后曾游到意大利岸边，随后遭当地人杀害而不得安葬。}}
